\documentclass{report}

\input{~/dev/latex/template/preamble.tex}
\input{~/dev/latex/template/macros.tex}

\title{\Huge{}}
\author{\huge{Nathan Warner}}
\date{\huge{}}
\pagestyle{fancy}
\fancyhf{}
\lhead{Warner \thepage}
\rhead{}
% \lhead{\leftmark}
\cfoot{\thepage}
%\setborder
% \usepackage[default]{sourcecodepro}
% \usepackage[T1]{fontenc}

\begin{document}
    % \maketitle
        \begin{titlepage}
       \begin{center}
           \vspace*{1cm}
    
           \textbf{Bash Mastery} \\
           The complete guide to BASH shell scripting
    
           \vspace{0.5cm}
            
                
           \vspace{1.5cm}
    
           \textbf{Nathan Warner}
    
           \vfill
                
                
           \vspace{0.8cm}
         
           \includegraphics[width=0.4\textwidth]{~/niu/seal.png}
                
           Computer Science \\
           Northern Illinois University\\
           August 3, 2023 \\
           United States\\
           
                
       \end{center}
    \end{titlepage}
    \tableofcontents
    \pagebreak \bigbreak \noindent
    \section{\LARGE Setting up scripts}
    \bigbreak \noindent 
    In order to write shell scripts, we must use the file extension .sh. When we begin the script, we must include a \textit{shebang}, which looks something like:
    \bigbreak \noindent 
    \begin{minted}{bash}
#!/usr/bin/bash
    \end{minted}
    \bigbreak \noindent 
    However, this could depend on the users system. To locate which file path to use, we can use the command:
    \begin{minted}{bash}
which bash
    \end{minted}
    \bigbreak \noindent 
    This will provide the path to use for the shebang.
    \bigbreak \noindent 
    The anatomy of a shell script can be described with the following parts:
    \bigbreak \noindent 
    \begin{enumerate}
        \item Shebang 
        \item Commands
        \item Exit statement (0=successful, 1-255=unsuccessful)
    \end{enumerate}
    \bigbreak \noindent 
    Example:
    \bigbreak \noindent 
    \begin{mdframed}[style=purplebox]
    \begin{minted}{bash}
#!/usr/bin/bash

echo "Hello World!"
exit 0
    \end{minted}
    \bigbreak \noindent 
    \end{mdframed}
    \bigbreak \noindent 
    \nt{The recommended file permissions for scripts is 744 (chmod 744 \textit{filename})}

    \bigbreak \noindent 
    \subsection{Adding scripts to PATH}
    \bigbreak \noindent 
    To add scripts to your PATH for BASH, we can open up our .bashrc, and add at the bottom of the file:
    \bigbreak \noindent 
    \begin{mdframed}[style=purplebox]
    \begin{minted}{bash}
export PATH="$PATH:$HOME/dirlocation"
    \end{minted}
    \end{mdframed}
    \bigbreak \noindent 
    This will append some directory to the end of our PATH variable. To do the same for the \textit{fish} shell, in the fish config file, we can add.
    \bigbreak \noindent 
    \begin{mdframed}[style=purplebox]
    \begin{minted}{bash}
set -gx PATH $PATH $HOME/somedirectory
    \end{minted}
    \end{mdframed}

    \pagebreak \bigbreak \noindent 
    \section{\LARGE Variables and Shell Expansions}
    \bigbreak \noindent 
    \subsection{User-Defined variables and parameter expansion}
    \bigbreak \noindent 
    \smallbreak \noindent
    \begin{definition}
        A \textbf{parameter} is any entity that stores values. In bash, we have three types:
        \begin{enumerate}
            \item Variables
            \item Positional Parameters
            \item Special Parameters
        \end{enumerate}
    \end{definition}
    \bigbreak \noindent 
    To define variables in our script, we can do:
    \begin{mdframed}[style=purplebox]
    \begin{minted}{bash}
identifier=value # NO WHITESPACE
name="nate" # Example
declare -i a=1 # Integer variable

# Parameter Expansion (Reference variables)
echo "Hello, ${name}!" 
    \end{minted}
    \bigbreak \noindent 
    \end{mdframed}

    \bigbreak \noindent 
    \subsection{Shell variables}
    \bigbreak \noindent 
    Shell variables are builtin variables that we can access but don't need to define ourself, some common shell variables are:
    \begin{itemize}
        \item PATH
        \item HOME 
        \item USER 
        \item HOSTNAME
        \item HOSTTYPE
    \end{itemize}
    \bigbreak \noindent 
    \subsection{Positional Parameters}
    \bigbreak \noindent 
    Positional parameters are variables that hold the command-line arguments to a script or function. They are denoted by numbers.
    \begin{enumerate}
        \item \$0 (Contains the name of the script)
        \item \$1, \$2 ... \$n (The first, second, third, etc. arguments to the script or function.)
        \item \$\# (The number of arguments passed to the script or function.)
        \item \$@ (All the arguments. When quoted ("\$@"), it treats each argument as a separate word. Useful for loops, more on this later)
        \item \$* (All the arguments. When quoted ("\$*"), it treats all arguments as a single word. Useful for loops, more on this later)
    \end{enumerate}

    \pagebreak \bigbreak \noindent 
    \subsection{Special Parameters}
    \bigbreak \noindent 
    These are variables that provide special functionality or information about the script or command's execution:
    \bigbreak \noindent 
    \begin{itemize}
        \item \$?: The exit status of the last executed command. \(0\) usually indicates success, and a non-zero value indicates an error.
        \item \$\$: The process ID (PID) of the currently executing script or shell instance.
        \item \$!: The process ID (PID) of the last backgrounded command.
        \item \$-: The current options set for the shell. For instance, if you used \texttt{set -x} for debugging, \texttt{x} would be part of the value.
        \item \$\_: The last argument of the previous command. Also sometimes used to get the last path argument to the \texttt{cd} command.
    \end{itemize}

    \bigbreak \noindent 
    \subsection{Parameter Expansion Tricks}
    \bigbreak \noindent 
    Default Values:
    \begin{itemize}
        \item \$\{parameter:-word\}: If parameter is unset or null, this expansion will return word. Otherwise, it returns the value of parameter.
        \item \$\{parameter:=word\}: If parameter is unset or null, it will be set to word.
    \end{itemize}
    \bigbreak \noindent 
    String Length:
    \begin{itemize}
        \item \${\#parameter}: Returns the length of the value of the parameter.
    \end{itemize}
    \bigbreak \noindent 
    Substring Expansion:
    \begin{itemize}
        \item \$\{parameter:offset:length\}: Extracts a substring from \$parameter starting at offset (0-indexed) and of length length.
    \end{itemize}
    \bigbreak \noindent 
    String Removal (Pattern Matching):
    \begin{itemize}
        \item \$\{parameter\#pattern\}: Removes the shortest match of pattern from the beginning of \$parameter.
        \item \$\{parameter\#\#pattern\}: Removes the longest match of pattern from the beginning of \$parameter.
        \item \$\{parameter\%pattern\}: Removes the shortest match of pattern from the end of \$parameter.
        \item \$\{parameter\%\%pattern\}: Removes the longest match of pattern from the end of \$parameter. 
    \end{itemize}
    \bigbreak \noindent 
    String Replacement:
    \begin{itemize}
        \item \$\{parameter/pattern/string\}: Replaces the first match of pattern with string in \$parameter.
        \item \$\{parameter//pattern/string\}: Replaces all matches of pattern with string in \$parameter
    \end{itemize}
    \bigbreak \noindent 
    Variable Indirection:
    \begin{itemize}
        \item \$\{!parameter\}: Treats the value of parameter as the name of another variable, and fetches the value of that variable.
    \end{itemize}
    \bigbreak \noindent 
    Case Modification:
    \begin{itemize}
        \item \$\{parameter$\land$\}: Capitalizes the first letter of the value.
        \item \$\{parameter$\land\land$\}: Capitalizes all letters of the value.
        \item \$\{parameter,\}: Converts the first letter to lowercase.
        \item \$\{parameter,,\}: Converts all letters to lowercase.
    \end{itemize}

    \pagebreak \bigbreak \noindent 
    \subsection{Command Substitution}
    \bigbreak \noindent 
    \begin{concept}
		 Command substitution can be used to:
	\end{concept}
    \begin{itemize}
        \item Save the output of commands in variables
        \item Use the output of one command \textit{inside} another command
    \end{itemize}
    \bigbreak \noindent 
    The syntax for this is:
    \begin{mdframed}[style=purplebox]
    \begin{minted}{bash}
$(command)
# Example...
time=$(date +%H:%M:%S)
echo "Hello, the current time is ${time}"
    \end{minted}
    \bigbreak \noindent
    \end{mdframed}

    \bigbreak \noindent 
    \subsection{Arithmetic Expansion}
    \bigbreak \noindent 
    The syntax for \textit{Arithmetic Expansion} is:
    \begin{mdframed}[style=purplebox]
    \begin{minted}{bash}
$((expression))
# Example...
echo $((1+1)) # 2
# When dealing with arithmetic expansion, we do not need a $ to reference variables
x=1
y=1
echo $(( x + y ))
    \end{minted}
    \bigbreak \noindent
    \end{mdframed}
    
    \bigbreak \noindent 
    \subsection{Dealing with floating point numbers}
    \bigbreak \noindent 
    To be able to do floating point arithmetic in our scripts, we need to use the \textbf{bc} command.
    \bigbreak \noindent 
    Example:
    \begin{mdframed}[style=purplebox]
    \begin{minted}{bash}
echo "scale=2; 5/2" | bc # 2.50
result=$(echo "scale=2; 5/2" | bc)
# scale sets the precision of the output
    \end{minted}
    \bigbreak \noindent
    \end{mdframed}
    \bigbreak \noindent 

    \pagebreak \bigbreak \noindent 
    \subsection{Tilde Expansion}
    \bigbreak \noindent 
    I'm sure you're already familiar with using tilde to jump to your home directory, but we can also use $\sim$- to jump between our current directory, and our home directory

    \bigbreak \noindent 
    \subsection{Brace Expansion}
    \bigbreak \noindent 
    We have two types of brace expansions:
    \begin{itemize}
        \item String lists
        \item Range lists
    \end{itemize}
    \bigbreak \noindent 
    Here is examples of what we can do with brace expansion:
    \bigbreak \noindent 
    \begin{mdframed}[style=purplebox]
    \begin{minted}{bash}
echo {jan,feb,march} # jan feb march NO WHITESPACE IN BRACES
echo {1..5} # 1 2 3 4 5
echo {1..10..2} # 1 3 5 7 9
echo {a..e} # a b c d e
echo {a,b}{1,2,3} # (Cartesian product...) a1 a2 a3 b1 b2 b3
# Useful for commands...
mkdir dir_{1..3}.txt
touch file_{1..5}.txt
    \end{minted}
    \bigbreak \noindent
    \end{mdframed}

    \pagebreak \bigbreak \noindent 
    \section{\LARGE How Bash Processes Commands}
    \bigbreak \noindent 
    Bash uses a 5 step process to interpret a command
    \bigbreak \noindent 
    \begin{enumerate}
        \item \textbf{Tokenisation:} A token is a sequence of characters that is considered as a single unit by the shell. The shell determines were a token starts and ends with the following special (meta) characters
            \begin{itemize}
                \item |
                \item \&
                \item ;
                \item ()
                \item < >
                \item Space, tab, newline
            \end{itemize}
            The shell then determines whether these tokens are words, or operators. A \textbf{word} is a token that does not contain an unquoted metacharacter. \textbf{Operators} are tokens that contain at least one unquoted metacharacter. This makes quoting a key concept in how the shell operates
        \item \textbf{Command identification:} The shell breaks the command up into either simple, or compound commands. \textbf{Simple commands} are just a bunch of individual words, and each simple command is terminated by a control operator. \textbf{Compound commands} provide bash with its programming constructs, such as if statements, for loops, while loops, etc...
        \item \textbf{Expansions:}
        \item \textbf{Quote removal:} We add quotes to control how the command is interpreted, so this step will simply remove all those supportive quotes.
        \item \textbf{Redirection:}
    \end{enumerate}
    After these 5 steps are completed, bash will then execute the command that is left over.

    \bigbreak \noindent 
    \subsection{Quoting}
    \bigbreak \noindent 
    \begin{concept}
		 Quoting is about \textbf{removing special meanings}. There are three types of quoting:
	\end{concept}
    \begin{itemize}
        \item \textbf{Backslash (\textbackslash):} This removes special meaning from next character
        \item \textbf{Single Quotes:} Removes special meaning from all characters inside
        \item \textbf{Double Quotes:} Removes special meaning from all inside except dollar signs (\$) and backticks (\ \`\ )
    \end{itemize}

    \pagebreak \bigbreak \noindent 
    \subsection{Step 1. Tokenisation}
    \bigbreak \noindent 
    We learned earlier that whether or not a token is interpreted as a word or an operator, depends on if there are any \textbf{unquoted metacharacters}. In bash we have two types of operators:
    \bigbreak \noindent 
    \begin{minipage}[t]{0.5\textwidth}
    \textbf{Control operators:}
    \begin{itemize}
        \item Newline: command separator, similar to the semicolon (;).
        \item | used to send the output of one command as the input to another command.
        \item || Used to execute the command following it only if the command preceding it fails 
        \item \& This is used to execute a command in the background
        \item \&\& Used to execute the command following it only if the command preceding it succeeds 
        \item ; Acts as a command separator.
        \item ;; Used in the context of a case statement in shell scripting. It signifies the end of an option within a case block.
        \item ;\& Also used in a case statement. After executing the associated block for a matched pattern, the control will flow to the block of the next pattern without testing.
        \item ;;\& Another operator used in a case statement. The control will test the next pattern after executing the block for the matched pattern.
        \item |\& This is shorthand for 2>\&1 |. It redirects both standard output (stdout) and standard error (stderr) of the command before the pipe to the command after the pipe.
        \item ( used to group commands and execute them in a subshell. 
        \item ) used to group commands and execute them in a subshell. 
    \end{itemize}
   \end{minipage}
   \begin{minipage}[t]{0.5\textwidth}
       \textbf{Redirection Operators}
       \begin{itemize}
           \item < Redirects input for a command from a file rather than from the keyboard.
            \item > Redirects the output of a command to a file
            \item << Provides multiple lines of input to a command 
            \item >> Similar to >, but instead of overwriting the file, it appends to the file.
            \item <\& Duplicates one input file descriptor to another, allowing for more advanced redirections.
            \item >\& Duplicates one output file descriptor to another.
            \item >| (Clobber): This is used in conjunction with the noclobber option in Bash (set -o noclobber). It allows you to forcefully overwrite a file when output redirection is used, even if noclobber is set.
            \item <<- Similar to <<, but it allows leading tabs (not spaces) to be ignored, 
            \item <>  Opens a file in read-write mode for a command. 
       \end{itemize}
   \end{minipage}

   \pagebreak \bigbreak \noindent 
   \subsection{Step 2. Command Identification (Simple commands)}
   \bigbreak \noindent 
   As we stated earlier, we have two types of commands, simple, and complex. Let's have a look at a simple command.
   \bigbreak \noindent 
   \begin{center}
       $\underbrace{\text{echo}}_{\parbox{2cm}{\centering \text{Command}\\ \text{name}}} \underbrace{1\ 2\ 3}_{\parbox{2cm}{individual arguments}}$
   \end{center}
   \bigbreak \noindent 
   All simple commands are terminated by a control operator, which we have discussed earlier. In this case, it is a newline
   \bigbreak \noindent 
   Thus, "echo" is identified as the command, and the rest is identified as the commands arguments, since there are no \textbf{control operators}

   \bigbreak \noindent 
  \subsection{Step 2. Command Identification (Simple commands)}
  \bigbreak \noindent 
  \begin{concept}
		 Compound commands are essentially bash's programming constructs. Each compound command begins with a \textbf{reserved word} and is terminated by a \textbf{reserved word}.
	\end{concept}

  \bigbreak \noindent 
  \subsection{Step 3. Expansions}
  \bigbreak \noindent 
  The shell goes through 4 stages of expansion.
  \begin{itemize}
      \item \textbf{Stage 1:} Brace expansion
        \item \textbf{Stage 2:}
            \begin{itemize}
                \item Parameter expansion
                \item Arithmetic expansion
                \item Command substation
                \item Tilde expansion
            \end{itemize}
        \item \textbf{Stage 3:} Word splitting
        \item \textbf{Stage 4:} Globbing
  \end{itemize}
  \bigbreak \noindent 
  An important thing to know is that expansions in earlier stages, happen before expansions in later stages. This means if we had the code:
  \begin{mdframed}[style=purplebox]
  \begin{minted}{bash}
x=10
echo {1..$x}
  \end{minted}
  \bigbreak \noindent
  \end{mdframed}
  \bigbreak \noindent 
  We would \textbf{not} get the expected results, this is because the brace expansion will take place \textbf{before} the parameter expansion.
  \bigbreak \noindent 
  Another important thing to know is that items in \textbf{stage 2} have the same precedence, thus, they will be preformed left to right. Similar to how multiplication and division works in the rules of PEMDAS.

  \pagebreak \bigbreak \noindent 
  \subsection{Word splitting}
  \bigbreak \noindent 
  \begin{concept}
     \textbf{Word splitting}  is a process the shell performs to split the result of some unquoted expansions into separate words. Word splitting can have some very significant effects on how your commands are interpreted
  \end{concept}
  \bigbreak \noindent 
  Word splitting is only performed on the results of unquoted 
  \begin{itemize}
      \item Parameter expansions
        \item Command substitutions
        \item Arithmetic expansions
  \end{itemize}
  \bigbreak \noindent 
  The characters used to split words are governed by the IFS (Internal Field Separator) variable.
  \begin{itemize}
      \item Space, tab, and newline
  \end{itemize}
  \bigbreak \noindent 
  Suppose we have:
  \begin{mdframed}[style=purplebox]
  \begin{minted}{bash}
numbers="1 2 3 4 5"
touch ${numbers}
  \end{minted}
  \bigbreak \noindent
  \end{mdframed}
  \bigbreak \noindent 
  We will get 5 \textbf{different} files, labeled 1-5. This is because the parameter expansion was \textbf{unquoted}. Instead, we can do:
  \begin{mdframed}[style=purplebox]
  \begin{minted}{bash}
numbers="1 2 3 4 5"
touch "${numbers}"
  \end{minted}
  \bigbreak \noindent
  \end{mdframed}
  \bigbreak \noindent 
  This will prevent word splitting, and treat "1 2 3 4 5" as a single word. Creating just \textbf{one} file.
  \bigbreak \noindent 
  So we shall follow one simple rule. If we want the output of a:
  \begin{itemize}
      \item Parameter expansion 
        \item Command substation
        \item Arithmetic expansion
  \end{itemize}
  \bigbreak \noindent 
  To be considered as a \textbf{single word}, we must \textbf{wrap that expansion in double quotes!}

  \pagebreak \bigbreak \noindent 
  \subsection{Globbing}
  \bigbreak \noindent 
  \begin{concept}
      \textbf{Globbing} is used as shorthand for listing the files that a command should operate on. Globbing is \textbf{only} performed on words (not operators)
  \end{concept}
  \bigbreak \noindent 
  Globbing patters are words that contain unquoted \textbf{Special pattern characters}
  \begin{itemize}
      \item * (used as a wildcard for any word)
        \item ? (used as a wildcard for any character, very strict in terms of length)
        \item \textlbrackdbl\ \  (Like question mark, but we can control the characters)
  \end{itemize}
  \begin{mdframed}[style=purplebox]
  \begin{minted}{bash}
ls *.txt # will list all .txt files
ls file?.txt # will list all .txt files named "file[somecharacter]"
# Note: Anything with 2 or more characters after "file" will not be listed
ls file[ab].txt # will list either filea.txt, fileb.txt, or nothing
  \end{minted}
  \bigbreak \noindent
  \end{mdframed}
  \bigbreak \noindent 
  We can also treat the brackets as character classes...
  \begin{mdframed}[style=purplebox]
  \begin{minted}{bash}
ls file[a-z].txt
  \end{minted}
  \bigbreak \noindent
  \end{mdframed}

  \bigbreak \noindent 
  \subsection{Quote removal}
  \bigbreak \noindent 
  During the quote removal stage, the shell removes ALL unquoted backslashes, single quote characters, and double quote characters that \textbf{did not} result from a shell expansion.
  \bigbreak \noindent 

  \bigbreak \noindent 
  \subsection{Redirection}
  \bigbreak \noindent 
  \textbf{Data Streams}
  \begin{itemize}
      \item Stream 0: Standard input (stdin)
      \item Stream 0: Standard output (stdout)
      \item Stream 0: Standard error (stderr)
  \end{itemize}
  \bigbreak \noindent 
  \textbf{Standard input} provides us with an alternative way of providing input to a command, aside from using command line arguments. Default source is keyboard
  \bigbreak \noindent 
  \textbf{Standard output} contains the data that is produced after a successful command execution. Default source is terminal 
  \bigbreak \noindent 
  \textbf{Standard error} contains all error messages and status messages  that a command produces . Default source is terminal
  \bigbreak \noindent 
  Using redirection is all about changing the sources or destinations of these streams. For example, if we use the \textit{cat} command, we are changing the input stream source from the keyboard to a file.
  \pagebreak \bigbreak \noindent 
  For redirection we can use the following operators:
  \begin{itemize}
      \item > Redirects standard output to a file, overwriting the file if it exists.
        \item >> Redirects standard output to a file, appending to the file if it exists.
        \item < Takes input from a file
        \item 2> Redirects standard error to a file, overwriting the file if it exists.
        \item 2>> Redirects standard error to a file, appending to the file if it exists.
        \item \&> or >\& Redirects both standard output and standard error to a file, overwriting the file if it exists.
        \item \&>> Redirects both standard output and standard error to a file, appending to the file if it exists.
        \item | (pipe) Takes the standard output of one command and uses it as the standard input for another command.
  \end{itemize}

  \pagebreak \bigbreak \noindent 
  \section{\LARGE Requesting user input}
  \bigbreak \noindent 
  \subsection{Positional Parameters}
  \bigbreak \noindent 
  Recall from 2.3, where we discussed the concept of \textbf{positional parameters}, in this section, we will discuss how we can use them in our scripts.
  \bigbreak \noindent 
  Suppose we had some script call
  \begin{mdframed}[style=purplebox]
  \begin{minted}{bash}
./1.sh nate $HOME
  \end{minted}
  \bigbreak \noindent
  \end{mdframed}
  \bigbreak \noindent 
  Here we are passing two arguments to this script. For these arguments to be relevant, let's take a look at the script...
  \begin{mdframed}[style=purplebox]
  \begin{minted}{bash}
echo "Hello, I am ${0}"
echo "My name is, ${1}"
echo "My home directory is ${2}"
# Hello, I am ./1
# My name is, nate
# My home directory is /home/datura
  \end{minted}
  \bigbreak \noindent
  \end{mdframed}
  \bigbreak \noindent 
  So you can see we are making use of some \textbf{positional parameters} to gain access to the things that we passed to the script we it was executed.
  \bigbreak \noindent 
  \nt{It's important we use the curly brace expansion syntax for double digit numbers or we might get unexpected results}

  \bigbreak \noindent 
  \subsection{Special Parameters}
  \bigbreak \noindent 
  Now that we have learned how to use \textbf{positional parameters} in our scripts. We will take a look at \textbf{Special parameters}, these parameters give us information about the script. The definitions for these were described in 2.4
  \bigbreak \noindent 
  \begin{mdframed}[style=purplebox]
  \begin{minted}{bash}
echo "Exit code: ${?}"
echo "PID: ${??}"
# Exit code: 0 
# PID: 0
  \end{minted}
  \bigbreak \noindent
  \end{mdframed}

  \pagebreak \bigbreak \noindent 
  \subsection{The read command}
  \bigbreak \noindent 
  We can use the \textbf{read} command in bash to get input from the user, for example:
  \begin{mdframed}[style=purplebox]
  \begin{minted}{bash}
read name age
# Or
read -p "Enter your name: " name
read -p "Enter your age " age
  \end{minted}
  \bigbreak \noindent
  \end{mdframed}
  \bigbreak \noindent 
  Then when we run this script, we can enter some data. This data will be stored in the variables "name" and "age". Notice that the second example lets us define a prompt.
  \bigbreak \noindent 
  Other options we have for the read command are:
    \begin{itemize}
        \item -t: timeout, specified in seconds
        \item -s: make input s.t it is not visible in the terminal 
    \end{itemize}

    \bigbreak \noindent 
    \subsection{The select command}
    \bigbreak \noindent 
    Using the \textbf{select} command, we can create menus for the user, and have the shell carry out different commands based on what is selected. 
    \begin{mdframed}[style=purplebox]
    \begin{minted}{bash}
PS3="Enter: "
select day in a b c d e f 
do 
echo $day
break
done
    \end{minted}
    \bigbreak \noindent
    \end{mdframed}
    \bigbreak \noindent 
    There are some key things to notice here. For one, we can define the \textbf{PS3} variable to give the user a custom prompt. Second, this select command will by default result in a continuous loop. Thus, we must use the \textbf{break} command to terminate it.

    \pagebreak \bigbreak \noindent 
    \section{\LARGE Logic}
    \bigbreak \noindent 
    \subsection{Chaining commands} 
    \bigbreak \noindent 
    In bash we have the following operators:
    \begin{itemize}
        \item \& (send command to background)
        \item ; (chain commands)
        \item \&\& (only execute next command if previous was successful)
        \item || (only execute next command if previous was unsuccessful)
    \end{itemize}

    \bigbreak \noindent 
    \subsection{Test commands and conditional operators}
    \bigbreak \noindent 
    \begin{concept}
       If a test is evaluated to be true, the test will return an exit status of 0. If a test is evaluated to be false, the test will return an exit status of 1. The \textbf{test} command can be used in bash to compare different pieces of information.
    \end{concept}
    \bigbreak \noindent 
    The syntax for the test command is a set of square brackets [\ \ ]. Its important that you have a space separating the contents of the square bracket and the actual square brackets otherwise you will get a syntax error.
    \bigbreak \noindent 
    \nt{The set of square brackets is the \textbf{short hand form}, the long hand form is using the \textbf{test} keyword}
    \bigbreak \noindent 
    Example:
    \bigbreak \noindent 
    In conjunction with the test command, we can use the following operators:
    \bigbreak \noindent 
    \begin{minipage}[t]{0.47\textwidth}
        \textbf{File Tests:}
        \begin{itemize}
            \item -e FILE: True if FILE exists.
            \item -f FILE: True if FILE exists and is a regular file.
            \item -d FILE: True if FILE exists and is a directory.
            \item -s FILE: True if FILE exists and has a size greater than zero.
            \item -r FILE: True if FILE exists and is readable.
            \item -w FILE: True if FILE exists and is writable.
            \item -x FILE: True if FILE exists and is executable.
            \item -h FILE or -L FILE: True if FILE exists and is a symbolic link.
        \end{itemize}
    \end{minipage}
    \begin{minipage}[t]{0.47\textwidth}
        \textbf{String tests:}
        \begin{itemize}
            \item -z STRING: True if STRING has a length of zero.
            \item -n STRING: True if STRING has a length greater than zero.
            \item STRING1 = STRING2: True if STRING1 and STRING2 are the same.
            \item STRING1 != STRING2: True if STRING1 and STRING2 are not the same.
        \end{itemize}
    \end{minipage}

    \pagebreak \bigbreak \noindent 
    \begin{minipage}[]{0.47\textwidth}
        \textbf{Arithmetic Comparisons:}
        \begin{itemize}
            \item -eq: Equal to.
            \item -ne: Not equal to.
            \item -lt: Less than.
            \item -le: Less than or equal to.
            \item -gt: Greater than.
            \item -ge: Greater than or equal to
        \end{itemize}
    
    \end{minipage}
    \begin{minipage}[]{0.47\textwidth}
        \textbf{Logical Operators:}
        \begin{itemize}
            \item ! (logical not)
            \item -a (logical and)
            \item -o (logical or)
        \end{itemize}
    \end{minipage}

    \bigbreak \noindent 
    Example:
    \begin{mdframed}[style=purplebox]
    \begin{minted}{bash}
# Long hand
if test 1 -le 2; then
    echo "1 is less than 2"
fi
# Short hand...
if [ 1 -le 2 ]; then 
    echo "1 is less than 2"
fi
    \end{minted}
    \bigbreak \noindent
    \end{mdframed}
    \bigbreak \noindent 
    Additionally, we have the \textbf{double bracket}. This is known as the extended test command, offers several advantages and additional capabilities over the traditional test command. It allows for:
    \bigbreak \noindent 
    \begin{itemize}
        \item Word splitting does not occur
        \item Regex matching using =$\sim$
        \item String comparison using == or !=
        \item Compound comparisons using \&\& (and) and || (or)
    \end{itemize}

    \bigbreak \noindent 
    \subsection{If statements}
    The syntax for if statements
    \bigbreak \noindent 
    \begin{mdframed}[style=purplebox]
    \begin{minted}{bash}
if [ condition ]; then 

elif [ condition ]; then

else 

fi
    \end{minted}
    \bigbreak \noindent
    \end{mdframed}

    \pagebreak \bigbreak \noindent 
    \subsection{Case statements}
    \bigbreak \noindent 
    Syntax:
    \begin{mdframed}[style=purplebox]
    \begin{minted}{bash}
case "${var}" in 
    value)
        statement
        ;;
    value)
        statement
        ;;
    \end{minted}
    \bigbreak \noindent
    \end{mdframed}
    \bigbreak \noindent 

    \pagebreak \bigbreak \noindent 
    \section{\LARGE Process options and read from files}
    \bigbreak \noindent 
    \subsection{While loops}
    \bigbreak \noindent 
    Syntax:
    \begin{mdframed}[style=purplebox]
    \begin{minted}{bash}
while [[ test ]]; do
    statements;
done
    \end{minted}
    \bigbreak \noindent
    \end{mdframed}
    \bigbreak \noindent 
    So for example:
    \begin{mdframed}[style=purplebox]
    \begin{minted}{bash}
read -p "Enter a number: " num
while [[ $num -lt 5 ]]; do
    echo $num
    num=$(($num+1))
done
    \end{minted}
    \bigbreak \noindent
    \end{mdframed}


    \bigbreak \noindent 
    \subsection{Handling command line options}
    \bigbreak \noindent 
    For this we can use the \textbf{getopts} command to get any options that the user runs with the script
    \bigbreak \noindent 
    Syntax:
    \begin{mdframed}[style=purplebox]
    \begin{minted}{bash}
getopts "ab" varname # Defines two possible options, 'a', and 'b', stores in varname
getopts 'a:b:' varname # Lets the options take in values. Stores in varname
echo $OPTARG # Built in variable to house the option value
    \end{minted}
    \bigbreak \noindent
    \end{mdframed}
    \bigbreak \noindent 
    In order to ensure that the getopts command gets ALL options, we must write the command as the test for a while loop
    \begin{mdframed}[style=purplebox]
    \begin{minted}{bash}
while getopts 'a:b:' varname; do 
    echo $OPTARG
done
# Run script
./scriptname -a 2 -b 3
    \end{minted}
    \bigbreak \noindent
    \end{mdframed}

    \pagebreak \bigbreak \noindent 
    \subsection{Iterating over files with read-while loop}
    \bigbreak \noindent 
    \begin{concept}
        Read-while loops are while loops that use the \textbf{read} command as their test command
    \end{concept}
    \begin{mdframed}[style=purplebox]
    \begin{minted}{bash}
while read line; do
    echo "$line"
done < "${1}"
    \end{minted}
    \bigbreak \noindent
    \end{mdframed}
    \bigbreak \noindent 
    We can also use a technique called \textbf{process substation} to redirect standard input from a process
    \begin{mdframed}[style=purplebox]
    \begin{minted}{bash}
while read line; do
    echo "$line"
done > >(ls $HOME)
    \end{minted}
    \bigbreak \noindent
    \end{mdframed}

    \pagebreak \bigbreak \noindent 
    \section{\LARGE Arrays and For loops}
    \bigbreak \noindent 

    





  




  




    

    





    
\end{document}
